% cpu/overheads.tex
% mainfile: ../perfbook.tex
% SPDX-License-Identifier: CC-BY-SA-3.0

\section{开销}
\label{sec:cpu:Overheads}
%
\epigraph{不要在不了解材料的情况下设计桥梁，也不要在不了解底层硬件的情况下设计底层软件。}
	 {佚名}

本节将展示上一节列出的性能障碍的实际\IXpl{overhead}。
不过，首先需要对硬件系统架构有一个粗略的了解，这是下一节的主题。

\subsection{硬件系统架构}
\label{sec:cpu:Hardware System Architecture}

\begin{figure}
\centering
\resizebox{3in}{!}{\includegraphics{cpu/SystemArch}}
\caption{系统硬件架构}
\label{fig:cpu:System Hardware Architecture}
\end{figure}

\Cref{fig:cpu:System Hardware Architecture}
展示了一个八核计算机系统的概略示意图。
每个芯片包含一对CPU核心，每个核心都有自己的cache，
以及一个允许这对CPU相互通信的互连。
系统互连则允许四个芯片之间以及与主存之间进行通信。

数据以``\IXpl{cache line}''为单位在该系统中流动，
cache line是大小固定、地址对齐的2的幂次方内存块，
通常在32到256字节之间。
当CPU从内存加载一个变量到其寄存器时，
它必须首先将包含该变量的cache line加载到其cache中，
cache可以被视为一个硬件哈希表（hash table）。
类似地，当CPU将一个值从其寄存器存储到内存时，
它也必须将包含该变量的cache line加载到其cache中，
但还必须确保没有其他CPU拥有该cache line的副本。

\begin{figure}
\centering
\resizebox{3in}{!}{\includegraphics{cpu/SimpleWrite}}
\caption{一次``简单''存储的生命周期}
\label{fig:cpu:Lifetime of a Simple Store}
\end{figure}

例如，假设CPU~1向一个变量\co{x}写入，而该变量的cache line
在CPU~6的cache中。
结合\cref{fig:cpu:Lifetime of a Simple Store}，
以下过程的简化视图由这些步骤组成，
其中每一行都将\cref{fig:cpu:System Hardware Architecture}
简化为只显示CPU~1和~6、它们的store buffer和本地cache，
以及用灰色矩形表示的连接两个CPU的系统互连。

\begin{enumerate}
\item	CPU~1检查其本地cache，未找到该cache line。
	因此它将写操作记录在其store buffer中，
	如\cref{fig:cpu:Lifetime of a Simple Store}的A行所示。
\item	对该cache line的请求被转发到CPU~0和~1的互连，
	互连检查CPU~1的本地cache，未找到该cache line。
	由于没有任何变化，系统状态仍如
	\cref{fig:cpu:Lifetime of a Simple Store}的A行所示。
\item	该请求被转发到系统互连，
	如\cref{fig:cpu:Lifetime of a Simple Store}的B行所示，
	系统互连向其他三个芯片查询，
	得知该cache line由包含CPU~6和~7的芯片持有。
\item	该请求被转发到CPU~6和~7的互连，
	互连检查两个CPU的cache，在CPU~6的cache中找到了该值，
	如\cref{fig:cpu:Lifetime of a Simple Store}的C行所示。
\item	CPU~6将cache line转发给其互连，同时将该cache line
	从其cache中清除，然后CPU~6和~7的互连将cache line
	转发到系统互连，
	如\cref{fig:cpu:Lifetime of a Simple Store}的D行所示。
\item	系统互连将cache line转发到CPU~0和~1的互连，
	如\cref{fig:cpu:Lifetime of a Simple Store}的E行所示。
\item	CPU~0和~1的互连将cache line转发给CPU~1，
	如\cref{fig:cpu:Lifetime of a Simple Store}的F行所示。
\item	该cache line被存入CPU~1的cache中，
	如\cref{fig:cpu:Lifetime of a Simple Store}的G行所示。
\item	CPU~1现在可以完成写操作，用之前记录在store buffer中的值
	更新新到达的cache line的相关部分，
	如\cref{fig:cpu:Lifetime of a Simple Store}的H行所示。
\end{enumerate}

\QuickQuizSeries{%
\QuickQuizB{
	这是一个\emph{简化}的事件序列？
	它怎么\emph{可能}更复杂？
}\QuickQuizAnswerB{
	这个序列忽略了许多可能的复杂情况，包括：

	\begin{enumerate}
	\item	其他CPU可能正在并发地尝试对同一cache line执行
		内存引用操作。
	\item	该cache line可能已被只读复制到多个CPU的cache中，
		在这种情况下，需要从这些cache中清除。
	\item	当请求到达时，CPU~6可能正在操作该cache line，
		此时CPU~6可能需要暂缓请求，直到自己的操作完成。
	\item	CPU~6可能已经将该cache line从其cache中驱逐
		（例如，为了给其他数据腾出空间），
		因此当请求到达时，该cache line正在返回内存的途中。
	\item	cache line中可能发生了可纠正的错误，
		需要在数据被使用之前的某个时间点进行纠正。
	\end{enumerate}

	由于这些考虑因素，生产级的cache一致性机制极其
	复杂~\cite{Hennessy95a,DavidECuller1999,MiloMKMartin2012scale,DanielJSorin2011MemModel}。
%
}\QuickQuizEndB
%
\QuickQuizE{
	为什么必须从CPU~6的cache中清除该cache line？
}\QuickQuizAnswerE{
	如果不从CPU~6的cache中清除该cache line，
	那么CPU~1和~6可能对cache line中同一组变量持有不同的值。
	这种不一致性会极大地复杂化并行软件，
	因此明智的硬件架构师会避免这种情况。
}\QuickQuizEndE
}

这个简化的序列只是一门称为\emph{cache一致性协议}的
学科的开端~\cite{Hennessy95a,DavidECuller1999,MiloMKMartin2012scale,DanielJSorin2011MemModel}，
该协议在\cref{chp:app:whymb:Why Memory Barriers?}中有更详细的讨论\@。
从一个简单的内存写操作所触发的事件序列可以看出，
一条指令可能引发大量的协议流量，
这会显著降低你的并行程序的性能。

幸运的是，如果某个变量在一段时间内被频繁读取但从未被更新，
该变量可以被复制到所有CPU的cache中。
这种复制使所有CPU都能极快地访问这个\emph{读多写少}的变量。
\Cref{chp:Deferred Processing}介绍了充分利用这一重要的
硬件读多写少优化的同步机制。

\subsection{操作开销}
\label{sec:cpu:Costs of Operations}

\begin{table}
%\rowcolors{1}{}{lightgray}
\renewcommand*{\arraystretch}{1.1}
\centering\small
\tcresizewidth{
\begin{tabular}
  {
    ll
    S[table-format = 9.1]
    S[table-format = 9.1]
    r
  }
	\toprule
	\multicolumn{2}{l}{Operation}
			& \multicolumn{1}{r}{Cost (ns)}
				   & {\parbox[b]{.7in}{\raggedleft Ratio\\(cost/clock)}}
					    & CPUs \\
	\midrule
	\multicolumn{2}{l}{Clock period}
			   &   0.5 &    1.0 &		\\
	\midrule
	\multicolumn{2}{l}{Same-CPU}
			   &       &        & 0		\\
		& CAS      &   7.0 &   14.6 &		\\
		& lock     &  15.4 &   32.3 &		\\
	\midrule
	\multicolumn{2}{l}{On-Core}
			   &       &        & 224	\\
		& Blind CAS&   7.2 &   15.2 &		\\
		& CAS	   &  18.0 &   37.7 & 		\\
        \midrule
	\multicolumn{2}{l}{Off-Core}
			   &       &	    & 1--27	\\
		& Blind CAS&  47.5 &   99.8 & 225--251	\\
		& CAS	   & 101.9 &  214.0 &		\\
        \midrule
	\multicolumn{2}{l}{Off-Socket}
			   &       &        & 28--111	\\
		& Blind CAS& 148.8 &  312.5 & 252--335	\\
		& CAS	   & 442.9 &  930.1 &		\\
        \midrule
	\multicolumn{2}{l}{Cross-Interconnect}
			   &       &        & 112--223	\\
		& Blind CAS& 336.6 &  706.8 & 336--447	\\
		& CAS	   & 944.8 & 1984.2 &		\\
	\midrule
	\multicolumn{2}{l}{Off-System}
				&		&	      & \\
		& Comms Fabric  &         5 000 &      10 500 & \\
		& Global Comms  &   195 000 000 & 409 500 000 & \\
	\bottomrule
\end{tabular}
}
\caption{8 路 Intel Xeon Platinum 8176 @ 2.10\,GHz 系统中 CPU 0 视角的同步机制开销}
\label{tab:cpu:CPU 0 View of Synchronization Mechanisms on 8-Socket System With Intel Xeon Platinum 8176 CPUs at 2.10GHz}
\end{table}

对并行程序重要的一些常见操作的开销显示在
\cref{tab:cpu:CPU 0 View of Synchronization Mechanisms on 8-Socket System With Intel Xeon Platinum 8176 CPUs at 2.10GHz}中。
该系统的时钟周期约为0.5\,ns。
虽然现代微处理器在一个时钟周期内退役多条指令并不罕见，
但这些操作的开销仍然在第三列（标记为``Ratio''）中以时钟周期为单位进行了归一化。
关于这张表首先要注意的是，许多比率的值非常大。

同CPU的\acrmf{cas}操作消耗约7纳秒，
这个持续时间是时钟周期的十倍以上。
CAS是一种原子操作，硬件将指定内存位置的内容与指定的``旧''值进行比较，
如果相等，则存储指定的``新''值，此时CAS操作成功。
如果不相等，则内存位置保持其（意外的）值，CAS操作失败。
该操作是原子的，因为硬件保证在比较和存储之间
内存位置不会被更改。
在x86上，CAS功能由\co{lock;cmpxchg}指令提供。

``同CPU''前缀意味着当前对给定变量执行CAS操作的CPU也是最后一个访问
该变量的CPU，因此相应的cache line已经在该CPU的cache中。
类似地，同CPU的lock操作（由一次锁获取和一次锁释放组成的``往返''对）
消耗超过15纳秒，即超过30个时钟周期。
lock操作比CAS更昂贵，因为它需要对lock数据结构执行两次原子操作，
一次用于获取，另一次用于释放。

核内操作涉及共享同一核心的硬件线程之间的交互，
其开销与同CPU操作大致相同。
这并不令人惊讶，因为这两个硬件线程也共享完整的cache层次结构。

在blind CAS的情况下，软件在不查看内存位置的情况下指定旧值。
这种方法适用于尝试获取锁的场景。
如果未锁定状态用零表示，锁定状态用值一表示，
那么对lock执行CAS操作，指定零为旧值、一为新值，
如果锁尚未被持有，就会获取锁。
关键在于只有一次对内存位置的访问，即CAS操作本身。

相比之下，普通CAS操作的旧值来自于某个先前的load。
例如，要实现原子递增，先将共享变量的当前值加载到一个机器寄存器中，
然后递增以在另一个机器寄存器中产生新值。
然后在CAS操作中，第一个寄存器被指定为旧值，第二个寄存器被指定为新值。
如果共享变量的值在此期间没有变化，
CAS操作将存储新值，从而递增该共享变量。
然而，如果共享变量的值已经改变，旧值将不匹配，
导致比较失败，从而使CAS操作失败，不再对该共享变量做进一步更改。
关键在于现在有两次对内存位置的访问：load和CAS\@。

因此，核内blind CAS只消耗约7纳秒，
而核内CAS消耗约18纳秒，这并不令人意外。
非blind情况下额外的load并非免费的。
也就是说，这些操作的开销与同CPU的CAS和lock分别相似。

\QuickQuiz{
	\Cref{tab:cpu:CPU 0 View of Synchronization Mechanisms on 8-Socket System With Intel Xeon Platinum 8176 CPUs at 2.10GHz}
	显示CPU~0与CPU~224共享一个核心。
	然而，CPU 0与CPU 1共享一个核心不是更合理吗？
}\QuickQuizAnswer{
	对这种看法表示理解是很容易的，但文件
	\path{/sys/devices/system/cpu/cpu0/cache/index0/shared_cpu_list}
	中确实包含字符串\co{0,224}。
	因此，CPU~0的超线程伙伴确实是CPU~224。
	有人推测这种编号方式能让简单的应用程序和调度器表现更好，
	理由是在许多工作负载上第二个超线程并不能提供大量额外性能。
	这种推测假设简单的应用程序和调度器会按数字顺序使用CPU，
	将较弱的超线程伙伴CPU留到所有核心都在使用时才启用。
}\QuickQuizEnd

跨核（不同核心但在同一socket上）的blind CAS消耗近50纳秒，
即近100个时钟周期。
用于此cache miss测量的代码在一对CPU之间来回传递cache line，
因此此cache miss是从另一个CPU的cache中满足的，而不是从内存。
非blind CAS操作（如前所述，需要查看变量的旧值并存储新值）
消耗超过100纳秒，即超过200个时钟周期。
想一想这意味着什么。
在执行\emph{一次}CAS操作所需的时间里，CPU本可以执行
超过\emph{两百}条普通指令。
这应该能说明不仅是细粒度锁的局限性，
而且是任何依赖细粒度全局一致性的同步机制的局限性。

如果这对CPU分别在不同的socket上，操作就会昂贵得多。
一次blind CAS操作消耗近150纳秒，即超过300个时钟周期。
一次普通CAS操作消耗超过400纳秒，即近\emph{一千}个时钟周期。

更糟糕的是，并非所有socket对都是平等的。
这个特定系统似乎是由一对四socket组件构成的，
当CPU位于不同组件中时会有额外的延迟惩罚。
在这种情况下，一次blind CAS操作消耗超过300纳秒，
即超过700个时钟周期。
一次CAS操作消耗近一微秒，即近2000个时钟周期。

\QuickQuizLabelRel{\QspeedOfLightAtoms}{1} % cann't put label inside QQSeries

\QuickQuizSeries{%
\QuickQuizB{
	硬件设计师难道不能被说服改善这种状况吗？
	为什么他们对这些单指令操作如此低劣的性能感到满意？
}\QuickQuizAnswerB{
	硬件设计师\emph{一直}在致力于解决这个问题，
	并且咨询了不亚于已故物理学家Stephen Hawking这样的权威人物。
	Hawking的观察是，硬件设计师有两个基本
	问题~\cite{BryanGardiner2007}：

	\begin{enumerate}
	\item	有限的光速，以及
	\item	物质的原子性质。
	\end{enumerate}

\begin{table}
%\rowcolors{1}{}{lightgray}
\renewcommand*{\arraystretch}{1.1}
\centering\small
\begin{tabular}
  {
    ll
    S[table-format = 9.1]
    S[table-format = 9.1]
  }
	\toprule
	\multicolumn{2}{l}{Operation}
			& \multicolumn{1}{r}{Cost (ns)}
			& {\parbox[b]{.7in}{\raggedleft Ratio\\(cost/clock)}} \\
	\midrule
	\multicolumn{2}{l}{Clock period}
			&           0.4	&           1.0 \\
        \midrule
	\multicolumn{2}{l}{Same-CPU}
			&		&		\\
	& CAS		&          12.2	&          33.8 \\
	& lock		&          25.6	&          71.2 \\
        \midrule
        \multicolumn{2}{l}{On-Core}
			&		&		\\
	& Blind CAS	&          12.9	&          35.8 \\
	& CAS		&           7.0	&          19.4 \\
	\midrule
        \multicolumn{2}{l}{Off-Core}
			&		&		\\
	& Blind CAS	&          31.2	&          86.6 \\
	& CAS		&          31.2	&          86.5 \\
	\midrule
	\multicolumn{2}{l}{Off-Socket}
			&		&		\\
	& Blind CAS	&          92.4	&         256.7 \\
	& CAS		&          95.9	&         266.4 \\
	\midrule
	\multicolumn{2}{l}{Off-System}
			&		&		\\
	& Comms Fabric	&       2 600   &       7 220   \\
	& Global Comms	& 195 000 000	& 542 000 000   \\
	\bottomrule
\end{tabular}
\caption{16 核 2.8\,GHz Intel X5550 (Nehalem) 系统的同步机制性能}
\label{tab:cpu:Performance of Synchronization Mechanisms on 16-CPU 2.8GHz Intel X5550 (Nehalem) System}
\end{table}

	第一个问题限制了原始速度，第二个限制了小型化，
	而小型化反过来又限制了频率。
	而且这还没有考虑功耗问题，功耗目前将生产频率
	限制在远低于10\,GHz。

	此外，
	\cref{tab:cpu:CPU 0 View of Synchronization Mechanisms on 8-Socket System With Intel Xeon Platinum 8176 CPUs at 2.10GHz}
	（位于第~\cpageref{tab:cpu:CPU 0 View of Synchronization Mechanisms on 8-Socket System With Intel Xeon Platinum 8176 CPUs at 2.10GHz}页）
	代表的是一个相当大的系统，拥有不少于448个硬件线程。
	更小的系统通常能实现更低的延迟，如
	\cref{tab:cpu:Performance of Synchronization Mechanisms on 16-CPU 2.8GHz Intel X5550 (Nehalem) System}所示，
	该表代表一个只有16个硬件线程的小得多的系统。
	\cref{tab:cpu:CPU 0 View of Synchronization Mechanisms on 8-Socket System With Intel Xeon Platinum 8176 CPUs at 2.10GHz}
	中到``Off-Core''行（含该行）为止的各行也提供了类似的视角。

\begin{table}
%\rowcolors{1}{}{lightgray}
\renewcommand*{\arraystretch}{1.1}
\centering\small
\tcresizewidth{
\begin{tabular}
  {
    ll
    S[table-format = 9.1]
    S[table-format = 9.1]
    r
  }
	\toprule
	\multicolumn{2}{l}{Operation}
		& \multicolumn{1}{r}{Cost (ns)}
			& {\parbox[b]{.7in}{\raggedleft Ratio\\(cost/clock)}}
			& CPUs \\
	\midrule
	\multicolumn{2}{l}{Clock period}
				     &   0.5 &    1.0 &			  \\
        \midrule
	\multicolumn{2}{l}{Same-CPU} &       &        &	0		  \\
	& CAS			     &   6.2 &   13.6 &			  \\
	& lock			     &  13.5 &   29.6 &			  \\
        \midrule
	\multicolumn{2}{l}{On-Core}  &       &        &	6		  \\
	& Blind CAS		     &   6.5 &   14.3 &			  \\
	& CAS			     &  16.2 &   35.6 &			  \\
        \midrule
	\multicolumn{2}{l}{Off-Core} &       &        &	1--5		  \\
	& Blind CAS		     &  22.2 &   48.8 & 7--11		  \\
	& CAS			     &  53.6 &  117.9 &			  \\
	\midrule
	\multicolumn{2}{l}{Off-System}&       &        &		  \\
	& Comms Fabric		      & 5 000 & 11 000 &		  \\
	& Global Comms		      & 195 000 000 & 429 000 000 &	  \\
	\bottomrule
\end{tabular}
}
\caption{12 核 Intel Core i7-8750H @ 2.20\,GHz 系统中 CPU 0 视角的同步机制开销}
\label{tab:cpu:CPU 0 View of Synchronization Mechanisms on 12-CPU Intel Core i7-8750H CPU @ 2.20GHz}
\end{table}

	此外，较新的小规模单socket系统，
	如我正在键入本文的这台笔记本电脑，
	也具有更合理的延迟，如
	\cref{tab:cpu:CPU 0 View of Synchronization Mechanisms on 12-CPU Intel Core i7-8750H CPU @ 2.20GHz}所示。

	或者换个角度看，1990年代中期的一台64-CPU系统的跨互连延迟
	超过5微秒，因此即使是
	\cref{tab:cpu:CPU 0 View of Synchronization Mechanisms on 8-Socket System With Intel Xeon Platinum 8176 CPUs at 2.10GHz}
	中所示的那个八socket 448硬件线程的庞然大物，
	也比其25年前的同类产品改进了五倍以上。

	将硬件线程集成到单个核心中以及将多个核心集成到单个芯片上
	已经大大改善了延迟，至少在单个核心或单个芯片的范围内是如此。
	整个系统的延迟也有所改善，但仅约两倍。
	不幸的是，在过去几年里，光速和物质的原子性质
	并没有发生太大变化~\cite{NoBugsHare2016CPUoperations}。
	因此，空间和时间局部性是并发软件的首要关注点，
	即使在相对较小的系统上运行也是如此。

	\Cref{sec:cpu:Hardware Free Lunch?}
	将探讨硬件设计师还能做些什么来缓解并行程序员的困境。
}\QuickQuizEndB
%
\QuickQuizE{
	\Cref{tab:cpu:Performance of Synchronization Mechanisms on 16-CPU 2.8GHz Intel X5550 (Nehalem) System}
	（在第~\cpageref{tab:cpu:Performance of Synchronization Mechanisms on 16-CPU 2.8GHz Intel X5550 (Nehalem) System}页
	\QuickQuizARef{\QspeedOfLightAtoms}的答案中）
	显示核内CAS比同CPU的CAS和核内blind CAS都快。
	这是怎么回事？
}\QuickQuizAnswerE{
	我\emph{确实}对获得的数据感到惊讶，并对其有效性进行了严格验证。
	我持续得到了相同的结果。
	一个可能解释该观察结果的理论是：
	核心中的两个线程能够重叠其访问，
	而单个CPU必须顺序执行所有操作。
	遗憾的是，似乎没有公开文档解释Intel X5550 (Nehalem)系统
	为什么会表现出这种行为。
}\QuickQuizEndE
}                 % End of \QuickQuizSeries

\begin{table}
\rowcolors{1}{}{lightgray}
\renewcommand*{\arraystretch}{1.1}
\centering\small
\begin{tabular}{lrrrrr}
	\toprule
	Level &  Scope & Line Size &   Sets & Ways &    Size \\
	\midrule
	L0    &   Core &        64 &     64 &    8 &     32K \\
	L1    &   Core &        64 &     64 &    8 &     32K \\
	L2    &   Core &        64 &   1024 &   16 &   1024K \\
	L3    & Socket &        64 & 57,344 &   11 & 39,424K \\
	\bottomrule
\end{tabular}
\caption{8 路 Intel Xeon Platinum 8176 @ 2.10\,GHz 系统的缓存结构}
\label{tab:cpu:Cache Geometry for 8-Socket System With Intel Xeon Platinum 8176 CPUs @ 2.10GHz}
\end{table}

不幸的是，核内和socket内通信的高速度并非免费的。
首先，一个核心内只有两个CPU，一个socket内只有56个，
而整个系统有448个。
其次，如
\cref{tab:cpu:Cache Geometry for 8-Socket System With Intel Xeon Platinum 8176 CPUs @ 2.10GHz}所示，
核内cache与socket内cache相比非常小，
而socket内cache与该系统配置的1.4\,TB内存相比又非常小。
第三，再次参考该表，cache被组织为硬件哈希表，
每个桶只能容纳有限数量的条目。
例如，L3 cache的原始大小（``Size''）近40\,MB，
但每个桶（``Line''）只能容纳11个内存块（``Ways''），
每块最多64字节（``Line Size''）。
这意味着只需12字节的内存（当然是精心选择的地址）
就能使这个40\,MB的cache溢出。
另一方面，同样精心选择的地址可能会充分利用整个40\,MB。

空间局部性显然极其重要，数据在内存中的分散也同样重要。

I/O操作甚至更加昂贵。
如``Comms Fabric''行所示，
高性能（且昂贵！）的通信架构，如InfiniBand或众多专有互连，
端到端往返延迟约为5微秒，
在此期间可能已经执行了超过\emph{一万}条指令。
基于标准的通信网络通常需要某种协议处理，
这进一步增加了\IX{latency}。
当然，地理距离也会增加延迟，
光在光纤中绕地球一圈的光速延迟约为195\emph{毫秒}，
即超过4亿个时钟周期，
如\cref{tab:cpu:CPU 0 View of Synchronization Mechanisms on 12-CPU Intel Core i7-8750H CPU @ 2.20GHz}的``Global Comms''行所示。

% Reference for Infiniband latency:
% http://www.hpcadvisorycouncil.com/events/2014/swiss-workshop/presos/Day_1/1_Mellanox.pdf
%     page 6/76 'Leading Interconnect, Leading Performance'
% Needs updating...

\QuickQuiz{
	这些数字大得离谱！
	我怎么可能理解它们？
}\QuickQuizAnswer{
	拿一卷卫生纸。
	在美国，每卷通常有大约350--500张。
	撕下一张代表一个时钟周期，放到一边。
	现在展开剩下的纸卷。

	展开的那堆卫生纸大致代表一次\IXacr{cas} cache miss。

	对于更昂贵的系统间通信延迟，
	用几卷（或几箱）卫生纸来代表通信延迟。

	重要安全提示：
	在挪用卫生纸时，务必考虑到与你同住的人的需求，
	特别是在2020年或类似的时期，
	那时商店货架上的卫生纸和很多其他东西都被抢购一空。

	此外，对于从事内核代码工作的人来说，
	CPU在cache miss期间禁用中断，
	类似于你在展开一卷卫生纸时屏住呼吸。
	\emph{你}能在屏住呼吸的同时展开多少卷卫生纸？
	你可能需要避免在那么多次cache miss期间禁用中断。\footnote{
		感谢Matthew Wilcox提出这个屏住呼吸的类比。}
}\QuickQuizEnd

\subsection{硬件优化}
\label{sec:cpu:Hardware Optimizations}

自然会问硬件在这方面提供了什么帮助，答案是``相当多！''

一种硬件优化是大容量cache line。
这可以提供显著的性能提升，特别是当软件顺序访问内存时。
例如，给定64字节的cache line和软件访问64位变量，
第一次访问由于光速延迟（如果不是其他原因的话）仍然会很慢，
但剩余的七次可以非常快。
然而，这种优化有一个阴暗面，即\IX{false sharing}，
当同一cache line中的不同变量被不同的CPU更新时就会发生，
导致高cache miss率。\footnote{
	这种情况有时被称为``cache thrashing''
	或``cacheline bouncing''。}
软件可以使用许多编译器中提供的对齐指令来避免false sharing，
添加这样的指令是调优并行软件的常见步骤。

第二种相关的硬件优化是cache预取，
硬件对连续访问做出反应，预取后续的cache line，
从而避免这些后续cache line的光速延迟。
当然，硬件必须使用简单的启发式方法来确定何时预取，
而这些启发式方法可能被许多应用程序复杂的数据访问模式所欺骗。
幸运的是，一些CPU系列为此提供了特殊的预取指令。
不幸的是，这些指令在一般情况下的有效性仍然存在争议。

第三种硬件优化是store buffer，
它允许一串store指令快速执行，
即使这些store指向非连续地址且所需的cache line
都不在CPU的cache中。
这种优化的阴暗面是内存乱序，
详见\cref{chp:Advanced Synchronization: Memory Ordering}。

第四种硬件优化是推测执行，
它可以让硬件充分利用store buffer而不导致内存乱序。
这种优化的阴暗面可能是能源效率低下和
当推测执行出错需要回滚重试时性能降低。
更糟糕的是，Spectre和Meltdown~\cite{JannHorn2018MeltdownSpectre}
的出现表明，硬件推测还可能使侧信道攻击成为可能，
这些攻击能击败内存保护硬件，
使无特权进程读取它们不应访问的内存。
显然，推测执行与云计算的结合需要相当多的改进！

有人可能会争辩说，如果人们能合理行事，
侧信道攻击的缓解措施就不是必需的。
然而，远程可访问的计算机系统确实经常受到
有组织犯罪和国家行为体的攻击，
更不用说无聊的青少年了。
有句老话说``少数人就能搞砸所有人的事''，
但现实是远程可访问的计算机系统必须被积极地防御攻击。

第五种硬件优化是大容量cache，
允许单个CPU在更大的数据集上操作而不产生昂贵的cache miss。
虽然大容量cache可能会降低\IXh{energy}{efficiency}和
\IXh{cache-miss}{latency}，
但生产微处理器上不断增长的cache大小证明了这种优化的强大。

最后一种硬件优化是读多写少的复制，
其中被频繁读取但很少更新的数据存在于所有CPU的cache中。
这种优化允许极其高效地访问读多写少的数据，
是\cref{chp:Deferred Processing}的主题。

\begin{figure}
\centering
\resizebox{3in}{!}{\includegraphics{cartoons/Data-chasing-light-wave}}
\caption{硬件与软件：并肩作战}
\ContributedBy{fig:cpu:Hardware and Software: On Same Side}{Melissa Broussard}
\end{figure}

简而言之，硬件工程师和软件工程师实际上是站在同一边的，
都在努力让计算机跑得更快，尽管物理定律竭尽全力加以阻挠，
如\cref{fig:cpu:Hardware and Software: On Same Side}
中所夸张描绘的那样，其中我们的数据流正竭尽全力超越光速，
还受到非零原子大小的进一步阻碍。
下一节将讨论硬件工程师可能（也可能无法）做的一些额外事项，
这取决于近期研究转化为实践的程度。
软件对这一崇高目标的贡献将在本书其余章节中概述。
